\section{完整证明}
\label{app:proofs}

\subsection{定理~\ref{thm:gd} 的证明（GD 线性收敛）}

设 $A = Q\Lambda Q^\top$，$0 < \mu \le \lambda_i \le L$。GD 更新 $x_{k+1} = x_k - \eta(A x_k - b)$，误差 $e_k = x_k - x^*$ 满足 $e_{k+1} = (I - \eta A) e_k$。在特征基下各模态衰减为 $|1 - \eta\lambda_i|$，故
\[
  \|e_{k+1}\|_2 \le \max_i |1 - \eta\lambda_i| \cdot \|e_k\|_2 = \rho(\eta)\,\|e_k\|_2.
\]
$\rho(\eta) = \max\{|1 - \eta\mu|, |1 - \eta L|\}$ 在 $\eta = 2/(\mu+L)$ 处取最小值 $\rho^* = (\kappa-1)/(\kappa+1)$。$\square$

\subsection{定理~\ref{thm:hb} 的证明（Polyak 最优 Heavy-ball）}

在特征基下，Heavy-ball 对每个模态 $\lambda$ 给出 $2\times 2$ 矩阵 $T(\lambda)$，特征方程 $z^2 - (1 - \eta\lambda + \beta) z + \beta = 0$。最优情形要求所有 $\lambda\in[\mu,L]$ 满足 $(1 - \eta\lambda + \beta)^2 \le 4\beta$，端点同时取等：
\[
  1 - \eta\mu + \beta = 2\sqrt\beta,\quad -(1 - \eta L + \beta) = 2\sqrt\beta.
\]
解得 $\beta^* = \bigl(\frac{\sqrt L - \sqrt\mu}{\sqrt L + \sqrt\mu}\bigr)^2$，$\rho^* = (\sqrt\kappa - 1)/(\sqrt\kappa + 1)$，$\eta^* = 4/(\sqrt L + \sqrt\mu)^2$。$\square$

\subsection{定理~\ref{thm:hb-cheb} 的证明（HB 是 Chebyshev 半迭代的定常极限）}

Chebyshev 半迭代（Hageman--Young 形式）：
\[
  \omega_1 = \frac{1}{1 - \sigma^2/2},\quad
  \omega_{k+1} = \frac{1}{1 - \sigma^2 \omega_k / 4},\quad
  x_{k+1} = \omega_{k+1}\left(\frac{r_k}{d} + x_k - x_{k-1}\right) + x_{k-1},
\]
其中 $d = (L+\mu)/2$，$\sigma = (L-\mu)/(L+\mu)$。令 $\omega_k \to \omega_\infty$，不动点方程给出 $\omega_\infty = 2/(1 + \sqrt{1-\sigma^2})$。可验证 $\eta_\infty = \omega_\infty/d = \eta^*$、$\beta_\infty = \omega_\infty - 1 = \beta^*$，故 Chebyshev 在系数收敛后退化为 Polyak 最优 Heavy-ball。实验 25 末段 gap $1.7\times 10^{-13}$ 与该解释一致。$\square$

\subsection{命题~\ref{prop:adam-jacobi} 的证明}

$\beta_1 = 0$ 时 $\hat m_k = g_k$；$v_k$ 趋于 $\mathbb{E}[g \odot g]$。对二次问题 $g(x) = A(x - x^*)$，当 $A$ 对角且 $\Sigma \propto I$ 时 $\mathbb{E}[g \odot g]\propto \mathrm{diag}(A)^2$，故 $P_k \approx \eta\,\mathrm{diag}(A)^{-1}$，恰为 Jacobi 预条件器。$\square$

\subsection{定理~\ref{thm:ns} 与命题~\ref{prop:basin} 的证明（Newton--Schulz 收敛盆）}

定理~\ref{thm:ns} 的误差递推 $E_{k+1}=-\tfrac14 E_k^2(3I-E_k)$ 的推导见正文 \S\ref{sec:ns-muon}。这里补全收敛盆 $(0,\sqrt3)$ 的严格证明。

\textbf{第一步：约简到标量映射。}\quad 迭代 \eqref{eq:ns} 只经 $X_k^\top X_k$ 作用于 $X_k$。设 $X_k=P_k\Sigma_k Q_k^\top$ 为 SVD，则 $X_k^\top X_k=Q_k\Sigma_k^2 Q_k^\top$，代入得
\[
  X_{k+1}=\tfrac12 P_k\Sigma_k(3I-\Sigma_k^2)Q_k^\top=P_k\,\varphi(\Sigma_k)\,Q_k^\top,
  \qquad \varphi(\sigma)=\tfrac12\sigma(3-\sigma^2).
\]
故奇异向量 $P_k\equiv P_0,\ Q_k\equiv Q_0$ 不随迭代改变，全部动力学集中在奇异值上：$\sigma^{(i)}_{k+1}=\varphi(\sigma^{(i)}_k)$，各分量解耦。$X_k\to U=P_0 Q_0^\top$（极分解正交因子）当且仅当所有 $\sigma^{(i)}_k\to1$。

\textbf{第二步：标量映射 $\varphi$ 在 $(0,\sqrt3)$ 上收敛到 $1$。}\quad
不动点方程 $\varphi(\sigma)=\sigma$ 化为 $\sigma(1-\sigma^2)=0$，正实轴上唯一非零不动点为 $\sigma^*=1$。导数 $\varphi'(\sigma)=\tfrac32(1-\sigma^2)$，于是 $\varphi'(1)=0$，$\sigma^*=1$ 为超吸引不动点。

\emph{(a) $\sigma\in(0,1]$。}\quad 此区间上 $\varphi'(\sigma)\ge0$，$\varphi$ 单调不减，且 $\varphi(0)=0,\varphi(1)=1$，故 $\varphi$ 自映 $(0,1]$。差式 $\varphi(\sigma)-\sigma=\tfrac12\sigma(1-\sigma^2)>0$（$\sigma\in(0,1)$），故迭代序列单调递增、上有界于 $1$，从而收敛，且极限为该区间唯一不动点 $1$。

\emph{(b) $\sigma\in(1,\sqrt3)$。}\quad 此区间上 $\varphi'(\sigma)<0$，$\varphi$ 单调递减，$\varphi(1)=1,\ \varphi(\sqrt3)=0$，故 $\varphi\bigl((1,\sqrt3)\bigr)=(0,1)$。任何此区间内的 $\sigma_0$ 一步后进入 $(0,1)$，回到情形 (a)，单调递增收敛到 $1$。

合并 (a)(b)：$(0,\sqrt3)$ 整体是 $\sigma^*=1$ 的吸引盆，且迭代恒保持在 $(0,\sqrt3)$ 内（不会逃逸）。

\textbf{第三步：二次收敛率。}\quad 令 $e_k=\sigma_k-1$。在 $\sigma=1$ 处 Taylor 展开 $\varphi(\sigma)=\varphi(1)+\varphi'(1)(\sigma-1)+\tfrac12\varphi''(1)(\sigma-1)^2+\tfrac16\varphi'''(1)(\sigma-1)^3$，其中 $\varphi'(1)=0,\ \varphi''(\sigma)=-3\sigma\Rightarrow\varphi''(1)=-3,\ \varphi'''(\sigma)=-3\Rightarrow\varphi'''(1)=-3$，得
\[
  e_{k+1}=-\tfrac32 e_k^2-\tfrac12 e_k^3,
  \qquad |e_{k+1}|\le\tfrac32|e_k|^2\ \ (|e_k|\le1).
\]
这与正文矩阵递推 $E_{k+1}=-\tfrac14E_k^2(3I-E_k)$ 在标量化（$E_k=\sigma_k^2-1$）下完全一致：$\sigma_k^2-1\approx2e_k$ 时两式同阶。

\textbf{第四步：发散侧。}\quad 若 $\sigma_0>\sqrt3$，则 $3-\sigma_0^2<0$，故 $\sigma_1=\tfrac12\sigma_0(3-\sigma_0^2)<0$；当 $|\sigma|$ 较大时 $\varphi(\sigma)\approx-\tfrac12\sigma^3$，$|\sigma_k|$ 以三次速度发散。这解释了实验 30 中 $\sigma_0>\sqrt3$ 的 $11$ 个初值全部发散。$\square$

\subsection{定理~\ref{thm:muon-tr} 的证明}

定理~\ref{thm:muon-tr} 的证明见正文 \S\ref{sec:ns-muon}，基于 von Neumann 迹不等式取等条件。
